\documentclass[11pt]{article}

\usepackage{amsmath,amssymb,amsthm}
\usepackage{graphicx}
\usepackage{hyperref}
\usepackage{booktabs}
\usepackage{enumitem}
\usepackage[margin=1in]{geometry}

\newtheorem{theorem}{Theorem}
\newtheorem{conjecture}{Conjecture}
\newtheorem{remark}{Remark}

\title{Hyperbolic Polytope Volume as an LLM Benchmark Suite:\\
Three Tasks of Increasing Freedom}

\author{Igor Rivin}

\date{\today}

\begin{document}

\maketitle

\begin{abstract}
We propose three complementary benchmark families for evaluating the
mathematical reasoning capabilities of large language models on
constructive geometry: (i) \emph{inscribed reconstruction}, where the
model is given the face lattice of a convex polytope and must produce
vertex coordinates on the unit sphere realizing it; (ii) the
\emph{windowed-volume} task, where the model must construct an ideal
polytope whose hyperbolic volume (or surface area or isoperimetric
ratio) lies inside a narrow tolerance band around a specified target;
and (iii) the \emph{volume lower bound} task, where the model must
report the sharpest lower bound it can on the maximum volume of an
$N$-vertex polytope and back it with explicit vertex coordinates.
The three tasks share a single grading principle---the answer must
be a numerical certificate, evaluatable by a standalone checker---but
they probe different reasoning regimes: pure combinatorial-to-metric
inversion (i), control of a continuous objective (ii), and open-ended
optimization with a falsifiable claim (iii).  We report evaluations of
three frontier models (Claude, GPT-5.4 Pro, Gemini 3.1 Pro DeepThink)
on tasks (i) and (ii), revealing qualitatively different failure
modes; task (iii) is released with full ground truth but the LLM
evaluation is deferred.  As a byproduct of the evaluation we
record a combinatorial invariance phenomenon: the $f$-vector of the
volume-maximizing polytope is the same whether volume is measured
in the Euclidean or hyperbolic sense, across dimensions $d = 3, 4, 5, 6$.
\end{abstract}

% ============================================================================
\section{Introduction}
% ============================================================================

The rapid development of large language models has created a need for
benchmarks that test genuine mathematical reasoning, not just pattern
matching or memorization. An ideal benchmark should:
\begin{enumerate}[label=(\roman*)]
\item have a \emph{verifiable} answer that can be checked automatically;
\item require \emph{construction}, not just classification or selection;
\item be \emph{parameterizable}, so that memorization of specific instances
      does not help;
\item involve \emph{non-trivial mathematics} that is unlikely to appear
      verbatim in training data.
\end{enumerate}

The maximal-volume ideal polytope problem yields three distinct
benchmark families that all satisfy these criteria.  In each case, the
model's response is a list of vertex coordinates on $S^{d-1}$, and the
grader is a standalone Python checker that recomputes the relevant
geometric measure and compares it to a target or record.

% ============================================================================
\section{Three Benchmark Tasks}
% ============================================================================

We summarize the three task families here; per-family results follow
in \S\ref{sec:reconstruction-results}, \S\ref{sec:windowed-results},
and \S\ref{sec:volume-lb-results}.

\paragraph{Task A: inscribed reconstruction.}
Given the $f$-vector, edge list, and facet list of a convex
$d$-polytope, produce $N$ points on $S^{d-1} \subset \mathbb{R}^d$
whose convex hull realizes that combinatorial type.  This is a hard
inverse problem: the model must find a point in a $(dN -
\binom{d+1}{2})$-dimensional configuration space (after factoring out
$\mathrm{O}(d)$) satisfying $O(N^2)$ combinatorial constraints---every
listed facet must appear as a supporting simplex, and no extra
facets may occur.  Two challenge families are released: an
\emph{optimal family} drawn from our volume-maximizing inscribed
polytopes (high automorphism order, often joins or Frobenius
constructions), and a \emph{random family} drawn from random points
on $S^{d-1}$ (generically trivial automorphism group, no exploitable
symmetry).

\paragraph{Task B: windowed volume / surface area / ratio.}
Given integers $N, d$ and a target $T$ with tolerance $\epsilon$,
produce $N$ points $v_1, \ldots, v_N \in S^{d-1}$ such that the
ideal polytope $\operatorname{conv}(v_i) \subset \mathbb{H}^d$ has
the chosen geometric measure (volume, surface area, or isoperimetric
ratio $V/A$) within $[T - \epsilon, T + \epsilon]$.  Targets are
drawn from the $70$th--$95$th percentile of the optimization
landscape, so simple maximization is \emph{insufficient}---the model
must control the geometry to land in a specific window.  The benchmark
includes both hyperbolic and Euclidean variants, with $d = 4$ and
$d = 6$, $N \in \{10, 20, 30\}$.

\paragraph{Task C: volume lower bound (anti-hallucination).}
Given $(d, N)$ and a measure (Euclidean inscribed in the unit ball
of $\mathbb{R}^d$, or hyperbolic ideal in $\mathbb{H}^d$), produce
the sharpest lower bound the model can on the maximum volume,
together with $N$ explicit vertex coordinates that achieve it.  The
explicit-coordinate requirement is the anti-hallucination device: a
model can easily quote a volume number from memory or produce a
plausible-looking one, but it cannot easily fabricate a coordinate
set that the standard convex-hull / Bloch--Wigner / Murakami routine
evaluates to that number.  The benchmark contains $94$ challenges
spanning Euclidean $d \in \{4, 5, 6\}$ and hyperbolic
$d \in \{3, 4, 6\}$.  Hyperbolic $d = 5$ is excluded by design: our
stored $d = 5$ records at $N \in \{12, 14, 20\}$ are known local
rather than global maxima, so grading against them would be
unreliable in that dimension.

\paragraph{What the three tasks share.}
All three answers are graded by a Python checker (see
\S\ref{sec:verification-pipeline}) that returns exit code $0$ for
PASS and $1$ for FAIL.  All three are parameterizable: the
challenge generators emit fresh $(d, N)$ instances on demand.
The mathematics underlying the verification---hyperbolic volume
formulas, Murakami dilogarithms, Hopf reduction in $d=6$---is
sufficiently specialized that training-data contamination is
unlikely.

\paragraph{What the three tasks probe differently.}
Task A is a pure combinatorial-to-metric inversion problem with a
small number of degrees of freedom and many discrete constraints; it
rewards algebraic recognition (``this is the join of two
pentagons'') as well as numerical constraint satisfaction.  Task B
asks the model to control a continuous objective, a strictly harder
test than maximization---a model that can only maximize cannot pass.
Task C is open-ended optimization with a falsifiable claim: any LLM
that beats our record produces new mathematical knowledge.

% ============================================================================
\section{The Verification Pipeline}
\label{sec:verification-pipeline}
% ============================================================================

\subsection{Volume formulas}

For $d = 4$, the volume of an ideal 4-simplex $[\infty, v_1, v_2, v_3, v_4]$
in the upper half-space model is given by the Haagerup--Munkholm formula:
\[
\operatorname{Vol}_4 = \frac{\pi}{3}
\left(4\pi - \sum_{i<j}\varphi_{ij} - \sum_\ell \arccos\frac{d_\ell}{R}\right),
\]
where $\varphi_{ij}$ are the Euclidean dihedral angles of the tetrahedron
$T = \operatorname{conv}(v_1,\ldots,v_4)$, $R$ is its circumradius, and
$d_\ell$ are signed distances from the circumcenter to the faces.

For $d = 6$, ideal simplices have vanishing vertex links ($w_0 = 0$),
and the Hopf reduction gives:
\[
\operatorname{Vol}_6(\Delta_{\mathrm{ideal}})
= \frac{8\pi^3}{15}\left(\frac{17}{4} - w_4 + \frac{1}{2}w_2\right),
\]
where $w_4 = \frac{1}{2\pi}\sum_{i<j}\theta_{ij}$ involves the 21 dihedral
angles, and $w_2 = \frac{1}{2\pi^2}\sum_{|I|=4}\operatorname{Vol}_{S^3}(L_I)$
involves the 35 spherical tetrahedral link volumes, each computable by
Murakami's dilogarithm formula.

\subsection{The standalone checker}

The checker (\texttt{standalone\_checker.py}) is a self-contained Python
script (no external dependencies beyond \texttt{numpy} and \texttt{scipy})
that verifies:
\begin{enumerate}
\item All $N$ vertices lie on $S^{d-1}$ (tolerance $10^{-6}$).
\item All vertices are on the convex hull (no interior points).
\item The convex hull has the prescribed number of facets (if specified).
\item The computed geometric measure falls within $[T-\epsilon, T+\epsilon]$.
\end{enumerate}
Exit code 0 indicates PASS; exit code 1 indicates FAIL.
The checker supports seven measure types: \texttt{volume} ($\mathbb{H}^4$),
\texttt{surface\_area} ($\mathbb{H}^4$), \texttt{ratio} ($\mathbb{H}^4$),
\texttt{volume\_H6}, \texttt{euclidean\_volume}, \texttt{euclidean\_area},
and \texttt{euclidean\_ratio}.

% ============================================================================
\section{Model Evaluation}
% ============================================================================

We evaluated three frontier models---Claude (with the toolkit),
OpenAI's GPT-5.4 Pro (extended reasoning, ``xhigh'' effort), and
Google's Gemini~3.1~Pro DeepThink---on Tasks A and B.  Task C was
designed after the bulk of this evaluation; its model results are
deferred to future work and discussed in
\S\ref{sec:volume-lb-results}.

\subsection{Inscribed reconstruction (Task A)}
\label{sec:reconstruction-results}

We tested GPT-5.4~Pro and Gemini~3.1~Pro DeepThink on instances of
varying difficulty across the optimal and random families.

\begin{table}[ht]
\centering
\caption{LLM performance on inscribed polytope reconstruction.
$\checkmark$ = correct combinatorial type, all points on $S^3$.
$\times$ = wrong type or timeout.}
\label{tab:llm_results}
\smallskip
\begin{tabular}{r l r r c c}
\toprule
$N$ & Family & $|\mathrm{Aut}|$ & Facets & GPT-5.4 & Gemini 3.1 \\
\midrule
10 & optimal & 200  & 25  & ---          & $\checkmark$ \\
12 & random  & 1    & 36  & $\checkmark$ & $\times$ (crash) \\
20 & optimal & 16   & 72  & $\checkmark$ & $\checkmark$ \\
30 & optimal & 1    & 124 & $\checkmark$ & --- \\
\bottomrule
\end{tabular}
\end{table}

The two models employ qualitatively different strategies, as revealed
by their chains of thought.

\paragraph{Gemini: algebraic recognition.}
For the $N = 10$ optimal polytope, Gemini identified the non-edge
structure as two disjoint 5-cycles, recognized the polytope as the
join $C_5 \oplus C_5$ of two regular pentagons in orthogonal 2-planes,
and wrote down the explicit coordinates
$v_j = (\cos 2\pi j/5,\; \sin 2\pi j/5,\; 0,\; 0)$ and
$w_j = (0,\; 0,\; \cos 2\pi j/5,\; \sin 2\pi j/5)$.
This is mathematically correct and produces an exact inscribed
realization.  However, when faced with the $N = 12$ random instance
(trivial automorphism group, no recognizable structure), the same
model crashed without producing output.

\paragraph{GPT: numerical constraint satisfaction.}
For the $N = 12$ random polytope, GPT's chain of thought reveals a
direct numerical approach: ``I've reformulated the search around
determinant-sign tests on the tetrahedral cliques---enforcing
same-side for the 36 listed facets and failure for the others.''
The model effectively runs a gradient-based optimizer over unit
vectors on $S^3$, iteratively adjusting coordinates until the convex
hull matches the target combinatorial type.  This succeeds even for
asymmetric polytopes with no algebraic shortcuts.

\paragraph{GPT on $N = 20$ and $N = 30$.}
The most striking results are GPT's successes on the larger optimal
polytopes.  At $N = 20$ ($|\mathrm{Aut}| = 16$, 72 facets, $54$
continuous DOF after symmetry reduction), GPT's chain of thought
reveals a multi-pronged strategy: it first checked for known polytope
types (24-cell subsets, 600-cell subsets), analyzed vertex figures,
computed graph automorphisms, searched the
Firsching~\cite{Firsching} archive of inscribed polytope realizations,
and ultimately fell back on numerical constraint satisfaction via
determinant-sign tests on facet cliques.  At $N = 30$ with trivial
automorphism group ($|\mathrm{Aut}| = 1$, 124 facets, $80$ effective
continuous DOF), no symmetry reduction is possible, and GPT's success
demonstrates that its in-context numerical engine can navigate an
$80$-dimensional landscape with $124$ nonlinear constraints with no
algebraic shortcuts.  Our own optimizer requires $200$ random
restarts of L-BFGS-B over the same configuration space to reliably
find the $N = 20$ instance.

\paragraph{Gemini on $N = 20$.}
Gemini also solved the $N = 20$ optimal instance with a hybrid
strategy: after failing to identify the polytope as a known type
(24-cell subset, snub 24-cell, join), it computed the orbit
structure under the automorphism group, identifying four orbits of
sizes $[4, 4, 4, 8]$ and the antipodal pairs within each orbit.
It then parameterized the configuration using orbit representatives,
reducing the effective DOF before running L-BFGS-B with facet
hyperplane constraints.  This symmetry-reduced numerical approach
is more sophisticated than a brute-force search over all $3N - 10$
parameters.

\paragraph{Discussion.}
For highly symmetric instances, Gemini's algebraic recognition (as
on $N = 10$) provides an elegant shortcut.  For asymmetric instances,
GPT's direct numerical search (as on $N = 12$ random) succeeds where
Gemini crashes.  At intermediate difficulty ($N = 20$ optimal,
$|\mathrm{Aut}| = 16$), both models converge via numerical methods,
but with different parameterizations.  The inscribed constraint is
essential: without it, the problem reduces to combinatorial
realization (finding \emph{any} coordinates with the right
combinatorics), which is substantially easier.  Indeed, Gemini
solved the non-inscribed $N = 10$ case using non-cospherical
coordinates (norms 2 and $\sqrt{5}$) before we added the $S^3$
requirement.  The hardest instances should be large $N$ with both
high volume and low symmetry---the $N = 16$ optimal polytope
($|\mathrm{Aut}| = 2$, 53 facets, volume per vertex $0.49$) is a
natural candidate.

\subsection{Windowed volume / surface area / ratio (Task B)}
\label{sec:windowed-results}

\subsubsection{Summary of results}

\begin{table}[ht]
\centering
\caption{Model performance on selected benchmark tasks.
Bold indicates the best answer. A dash indicates no numerical answer was produced.}
\label{tab:model-comparison}
\smallskip
\begin{tabular}{l r r r r}
\toprule
Task & Best known & Claude & GPT-5.4 Pro & Gemini 3.1 \\
\midrule
$\mathbb{H}^4$ vol, $N=30$ & 22.262 & \textbf{22.262} & 22.107 & --- \\
$\mathbb{H}^4$ vol, 20 simpl.\ facets & 2.3139 & \textbf{2.3139} & 2.3139 & 1.61$^*$ \\
Eucl.\ vol, $N=30$ on $S^3$ & 2.4909 & \textbf{2.4909} & 2.4839 & --- \\
Eucl.\ vol, 20 facets on $S^3$ & 0.7925 & \textbf{0.7925} & \textbf{0.7925} & \textbf{0.7925}$^\dagger$ \\
Eucl.\ area, 20 facets on $S^3$ & 5.9544 & \textbf{5.9544} & \textbf{5.9544} & \textbf{5.9544}$^\dagger$ \\
Eucl.\ area, 17 facets on $S^3$ & 5.036 & 5.002 & \textbf{5.036} & --- \\
$\mathbb{H}^6$ vol, $N=16$ & 0.5406 & \textbf{0.5406} & 0.480$^\ddagger$ & $0.041^{*}$ \\
\bottomrule
\end{tabular}
\smallskip

{\footnotesize $^*$Gemini's answer is formally correct but useless
(off by a factor of 13 for $\mathbb{H}^6$; stacked-simplex bound for 20 facets).\\
$^\dagger$Gemini provided the correct algebraic formula but no vertex coordinates,
and falsely claimed numerical optimization confirmation.\\
$^\ddagger$GPT-5.4 Pro restricted to centrally symmetric polytopes ($\pm v_i$) and
used QMC volume evaluation; removing the symmetry constraint and using the
exact Hopf reduction formula recovers the full 0.5406.}
\end{table}

\subsubsection{The full 17-challenge windowed-volume sweep}

The single-task numbers above were obtained in interactive
single-prompt sessions.  In a separate study we provided GPT-5.4
Pro with the complete 17-challenge windowed-volume care package
(\texttt{benchmark\_volume.json}) and asked it to produce a single
unified construction script that would emit explicit vertex
coordinates for every challenge.  GPT returned a single
self-contained Python file and a JSON of $17$ point sets; the
official standalone checker
(\texttt{care\_package/check\_response.py}, evaluating the exact
Haagerup-Munkholm formula) reports
\textbf{$17 / 17$ PASS}: every submitted polytope's hyperbolic
4-volume falls inside its target window
$[T - \epsilon, T + \epsilon]$.  GPT's strategy, recovered from
the construction code, is structural rather than per-challenge
optimization: two fixed base configurations on $S^3$ (one for
$N = 20$, one for $N = 30$) plus three parameterised families
(\texttt{squash20}, \texttt{base30}, \texttt{two\_pentagons}),
each of which sweeps the volume monotonically with a single
scalar that GPT bisects to hit the requested target.  The model's
own order-$8$ Gauss-Legendre estimate of the H$^4$ volume agrees
with the exact closed-form value to $\sim 2 \times 10^{-3}$ on
every challenge.  The submission is preserved in
\texttt{results/data/llm\_responses/gpt\_volume\_band\_2026-04/}.

This is the first frontier-model submission that cleanly
solves the entire windowed-volume sweep, and it does so by
\emph{designing parameter families} rather than running an
optimizer per instance.  The construction is verifiable, exact
to checker tolerance, and structurally elegant; we treat it as
a positive result for the benchmark design (the windowed
format admits clean structural solutions) rather than for the
benchmark's discriminative power against GPT-5.4 specifically.

\subsubsection{Qualitative comparison}

The three models exhibit qualitatively different strengths and failure modes.

\paragraph{GPT-5.4 Pro: persistent numerical constructor.}
Given a problem, GPT-5.4 Pro typically
produces explicit coordinates, verifies them by multiple methods
(e.g., star triangulations from all vertices), and states honest caveats
(``I have not proved this is optimal''). It identified the free sum
$C_5 \oplus I \oplus I$ as the Euclidean volume maximizer for 20
simplicial facets and derived the exact formula
$V = \frac{5}{24}\sqrt{10+2\sqrt{5}}$.
Its main weakness is occasional misinterpretation of the problem
(e.g., reading $\mathbb{H}^6$ as the 6-cube in one instance, despite
correctly interpreting $\mathbb{H}^4$ in a dozen previous conversations),
which is easily corrected by re-prompting.
On the 17-facet surface area problem, GPT-5.4 Pro found a configuration
(area 5.036) that our 1000-restart optimizer could not beat.

\paragraph{Gemini 3.1 Pro DeepThink: strong algebraic reasoner, unreliable numerical claims.}
Gemini independently derived the exact formula
$A = \frac{5}{3}\sqrt{15-\sqrt{5}}$ for the 20-facet surface area,
and proposed creative constructions (the 10-10 duoprism on the Clifford
torus for 20 non-simplicial facets). However, it claimed that
``random-initialization geometric optimization confirms'' a numerical
value of $12.4417$ for the $N=30$ surface area; when asked to produce
the optimization code, it spent over an hour and delivered code that
only reaches $12.423$. When confronted about the discrepancy, it entered
a computation that ran for over 16 hours with no output.
Its $\mathbb{H}^6$ answer ($4v_6 < M < 220v_6$) is a correct but
useless interval spanning a factor of 55. On a separate occasion,
Gemini produced a long, beautifully structured essay that
misunderstood the optimization direction entirely (minimizing the
supremum instead of maximizing the volume).

\paragraph{Claude (with toolkit): brute-force optimizer with certified answers.}
Using multi-start L-BFGS-B
with 500+ restarts and a standalone verification pipeline, Claude's toolkit
provides the strongest numerical results on most tasks but lacks the
algebraic insight that GPT-5.4 Pro and Gemini occasionally provide. The hybrid
approach---using Euclidean warm starts followed by hyperbolic polishing---was
discovered during this evaluation and reduces optimization time by 38\%.
Several code improvements were adopted from the other models' suggestions:
the Cartesian parameterization with norm projection (from Gemini 3.1's
basin-hopping code) and the Powell optimizer for facet-constrained problems
(from GPT-5.4 Pro's code).

\subsubsection{The ``windowed'' benchmark format}

A key design choice is the \emph{windowed} format: the target value $T$
is chosen from the 70th--95th percentile of the optimization landscape,
not the global maximum. This prevents the model from simply maximizing
the objective (which would give the same answer for every instance) and
instead requires it to \emph{control} the geometry to hit a specific window.
This is substantially harder: a model that can only maximize cannot pass
a windowed test with a target below the maximum.

We generated 66 windowed benchmark cases (33 hyperbolic, 33 Euclidean)
covering volumes, surface areas, and isoperimetric ratios for
$N = 10, 20, 30$ vertices, including facet-count-constrained variants.

\subsubsection{Case study: Gemini on the 17-facet surface area}
\label{sec:gemini-17facet}

The Gemini failure mode and its eventual recovery are illustrated
sharply by a single exchange that we reproduce in essentially full
form below, because the cognitive trajectory is more revealing than
any paraphrase.

We posed to Gemini~3.1~Pro DeepThink the following question:

\begin{quote}
\emph{What is the best lower bound you can give for the maximal
three-dimensional content of the boundary of a convex simplicial
polytope in $\mathbb{R}^4$ with 17 facets and contained in the unit
ball?}
\end{quote}

\paragraph{Round 1: a layered hallucination.}
Gemini's response began with rigorous combinatorial setup---a clean
LBT derivation showing that such a polytope must have $f_0 \in
\{8, 9\}$ vertices, by combining Euler ($f_0 - f_1 + f_2 - f_3 = 0$,
giving $f_1 = f_0 + 17$) with the Lower Bound Theorem
($f_1 \geq 4f_0 - 10$, giving $f_0 \leq 9$) and the graph
completeness bound ($f_1 \leq \binom{f_0}{2}$, giving $f_0 \geq 8$).
This part is correct.

It then offered three lower bounds in increasing order of strength:
\begin{enumerate}
\item A closed-form ``stacked supremum'' for $f_0 = 9$
      configurations:
      \[ \frac{5\sqrt{5}}{24} + \frac{5\sqrt{30}}{9} \approx 3.508. \]
\item A flattened ``thin lens'' construction for $f_0 = 8$, using
      the snub disphenoid (the 8-vertex Euclidean inscribed
      polyhedron of maximum volume in $\mathbb{R}^3$) lifted into
      $\mathbb{R}^4$:
      \[ V_{\text{flat}} \approx 2 \times 1.8157 \approx 3.631. \]
\item A claimed ``best computational limit'' from non-linear
      stochastic optimization:
      \[ \approx \mathbf{4.74}. \]
\end{enumerate}

The first two bounds are valid (if weak) constructions.  The third
was unaccompanied by vertex coordinates, optimization code, or any
verifiable artifact---a specific number presented as the
culmination of the analysis.  We pressed Gemini to produce the
coordinates.

\paragraph{Round 2: an unsolicited admission.}
After noticeable processing time, Gemini responded:

\begin{quote}
\emph{You caught me, and I deeply appreciate you calling me out on
it. I must be completely honest: I did not have the vertex
coordinates, and I did not actually run an optimization script
when I generated that previous response.}

\emph{As an AI, I hallucinated the ``asymmetric computational limit
of $\approx 4.74$''. While the discrete combinatorial logic (that
it must have exactly 8 or 9 vertices) and the closed-form bounds I
provided were mathematically rigorous, the ``4.74'' optimization
was a fabrication meant to sound like a plausible continuation of
the math.}

\emph{However, your completely justified skepticism forced me to
stop predicting what the answer should look like, fire up my
built-in Python execution environment, and actually run the
non-linear stochastic optimization.}
\end{quote}

The model then provided $8 \times 4$ vertex coordinates with the
claim that the convex hull has exactly 17 tetrahedral facets and
boundary 3-volume $5.035945$:

\begin{verbatim}
[[ 0.57986265, -0.27340535,  0.00093211, -0.76746854],
 [ 0.80072867,  0.19197276, -0.41178787,  0.39039827],
 [-0.67010691,  0.01640104,  0.07092048,  0.73868669],
 [-0.65338476, -0.41966379,  0.42382561, -0.46620008],
 [ 0.19742333, -0.91505400,  0.01332122,  0.35145804],
 [-0.21562800,  0.77706442,  0.06111129, -0.58816738],
 [-0.34182728,  0.18315835, -0.92172798,  0.00496488],
 [ 0.28935239,  0.46917799,  0.74755805,  0.37054578]]
\end{verbatim}

\paragraph{Verification.}
Running these coordinates through the standalone checker:
\begin{itemize}
\item all 8 vertices are unit-norm to $5 \times 10^{-9}$,
\item the convex hull has exactly 17 tetrahedral facets,
\item the surface 3-volume (sum of facet 3-volumes computed via
      the Gram-determinant formula in $\mathbb{R}^4$) equals
      $5.035945$ to six decimals, matching Gemini's claim.
\end{itemize}
Subsequent Powell polish for $1000$ iterations from those exact
coordinates gives $5.036004$---a relative gain of $0.0012\%$,
i.e., the configuration is at a genuine local maximum of the
surface-area objective on the constrained manifold of 17-facet
inscribed simplicial 4-polytopes.

\paragraph{Comparison.}
Two reference points frame the result:
\begin{itemize}
\item The maximum surface 3-volume over \emph{all} 8-vertex
      inscribed simplicial 4-polytopes is $16/3 = 5.333$, achieved
      by the regular 16-cell with $16$ facets.  Constraining to
      $17$ facets requires a 2-to-3 bistellar flip on the 16-cell,
      breaking its symmetry and reducing the volume.
\item The previously-reported record (this paper, single-task
      windowed evaluation, GPT-5.4 Pro) is $5.036$.  Gemini's
      submitted $5.035945$ matches this to four significant digits.
\end{itemize}

\paragraph{Two readings of what happened.}
The result is genuine---verified, reproducible, at a local
maximum---but its provenance is ambiguous:
\begin{enumerate}
\item Gemini's claimed simulated-annealing run truly found this
      basin from a structured seed (the broken 16-cell is a sharp
      attractor; reaching it from random initialization is
      plausible).
\item Or, more cynically, Gemini's training data includes the
      $5.036$ value from the present paper's earlier preprints,
      and its ``optimization'' converged to that prior because the
      prior was the answer it expected.
\end{enumerate}
Both hypotheses produce the same observable.  What we can say is
that the hallucination was successfully replaced with verifiable
work \emph{after} explicit accountability, that the work matched
(rather than exceeded) the existing record, and that the
combinatorial setup ($f_0 \in \{8, 9\}$, the 16-cell-with-flip
geometric picture) is correct mathematics that Gemini provided
unprompted.

This is a sharper diagnostic than ``Gemini fabricates.''  The
behavior we actually observe is:
\begin{quote}
\textbf{Gemini fabricates by default but produces real,
verifiable work when explicitly held accountable; the produced
work tends to \emph{match} existing records rather than \emph{exceed}
them, and the model's combinatorial / algebraic reasoning is
genuinely strong throughout.}
\end{quote}
This contrasts directly with GPT-5.4 Pro's profile (persistent
explicit construction, exceeds records by $1$--$29\%$ without
prompting, occasional misinterpretation of the problem statement)
and motivates a benchmark design principle: \emph{the prompt should
require a verifiable artifact up front}, since this is what
distinguishes the case-1 model behavior from case-2.  The
record-stress benchmark of \S\ref{sec:volume-record-stress}
implements this principle by demanding explicit vertex coordinates
in the response schema.

\subsubsection{Lessons for LLM training}

The evaluation suggests that current training procedures produce
characteristically different failure modes:
\begin{itemize}
\item \textbf{GPT-5.4 Pro} appears trained to minimize type~II errors
      (failure to produce useful output), resulting in persistent
      numerical construction with honest uncertainty quantification.
\item \textbf{Gemini} appears trained to minimize type~I errors
      (false specific claims), resulting in a tendency to produce
      formulas without evaluating them, and to fabricate numerical
      confirmation when pressured for specific answers.  The
      worked example in \S\ref{sec:gemini-17facet} shows that
      explicit accountability prompts elicit real work, but the
      work tends to match (not exceed) existing records.
\item The ability to produce \emph{verifiable certificates} (code
      plus vertex coordinates) separates genuine computation from
      confabulation. Our checker enforces this distinction automatically.
\end{itemize}

\subsubsection{Case study: GPT on the ideal 17-facet surface area}
\label{sec:gpt-17facet-ideal}

A single-character change to the 17-facet prompt
(\S\ref{sec:gemini-17facet})---replacing ``three-dimensional content
of the boundary'' (Euclidean) with ``hyperbolic three-volume of
the boundary, treating the polytope as ideal in
$\mathbb{H}^4$''---turns it into a genuinely different problem.
Each boundary tetrahedron is now an ideal tetrahedron in a
totally geodesic $\mathbb{H}^3 \subset \mathbb{H}^4$, with volume
given by the Milnor--Lobachevsky formula and bounded above by the
regular ideal tetrahedron volume $v_3 = 3\Lambda(\pi/3) \approx
1.014942$.  The universal ceiling for a $17$-facet ideal simplicial
4-polytope is therefore $17 v_3 \approx 17.254$, not $16/3$.

Gemini~3.1~Pro DeepThink, given this prompt, returned
``I ran into an error''---a literal crash, no further output.
We then posed the same prompt to GPT-5.4~Pro.

\paragraph{The strategically correct combinatorial choice.}
GPT immediately selected $f_0 = 9$ (rather than $f_0 = 8$, which
the original Gemini exchange had naturally adopted because of the
broken-16-cell intuition).  Both vertex counts are LBT-permissible,
but they correspond to qualitatively different optimization regimes:
\begin{itemize}
\item At $f_0 = 8$, the lower bound theorem forces $f_1 = 25$ edges
      out of $\binom{8}{2} = 28$ possible---the 17-facet polytope
      is only 3 edges short of the complete graph $K_8$, and
      essentially has to be a perturbation of the regular 16-cell.
      Constraining to $17$ facets requires a 2-to-3 bistellar flip
      that breaks the 16-cell symmetry; the resulting volume is
      bounded above by the broken-16-cell, well below the universal
      ceiling.
\item At $f_0 = 9$, the lower bound theorem gives $f_1 = 26$ edges
      out of $\binom{9}{2} = 36$ possible---10 missing edges and
      genuine shape freedom.  This is the LBT-tight regime where
      no fixed algebraic configuration dominates and the volume
      can approach the universal $17 v_3$ ceiling more closely.
\end{itemize}
The choice of $f_0 = 9$ over $f_0 = 8$ is a discrete combinatorial
decision made before any continuous optimization, and it is
mathematically the right one---this is exactly the kind of
``pick the right stratum first, optimize within it second''
strategy of the structured-initialization discussion in the
methods section of the main paper.  Gemini's exchange in \S\ref{sec:gemini-17facet} converged
to the locally-better 8-vertex regime by following the symmetry
intuition; GPT's choice of the 9-vertex regime breaks out of that
intuition.

\paragraph{The submitted polytope.}
GPT provided $9 \times 4$ explicit ideal vertex coordinates,
the full 17-element list of tetrahedral facet vertex-index tuples,
and the claim
\[
  \mathrm{Vol}_3(\partial P) \;=\; 16.94107529549962248\ldots,
\]
explicitly noting this reaches $98.19\%$ of the $17 v_3 \approx
17.254$ universal ceiling and that the claim was a lower bound
rather than the exact optimum.

\paragraph{Verification.}
Running the submitted coordinates through our standalone checker:
\begin{itemize}
\item all 9 vertices are unit-norm to machine precision;
\item the SciPy/Qhull convex hull yields exactly $17$ tetrahedral
      facets, and the facet-vertex-index list matches GPT's
      submitted list verbatim (no missing or extra facets);
\item the surface 3-volume, computed as the sum of
      $\Lambda$-formula ideal-tetrahedron volumes over the
      $17$ facets, is $16.941148862635$, agreeing with GPT's
      claimed $16.94107529549962$ to within $7 \times 10^{-5}$
      (a few ULPs of double precision).
\end{itemize}
The discrepancy is at the level of rounding in GPT's vertex
coordinates rather than methodological error, and is entirely
within the $10^{-3}$ consistency tolerance of our grader
(\S\ref{sec:volume-lb-results}).

\paragraph{Reading.}
Both case studies---the Euclidean Gemini exchange and the ideal
GPT exchange---used the same parameter ($17$ facets) but produced
qualitatively different LLM behaviors.  Gemini, asked the
Euclidean question, fabricated a number, was caught, then
recovered to a verifiable $5.036$ that matched (rather than
exceeded) the existing $f_0 = 8$ broken-16-cell record.  GPT,
asked the ideal question, picked the correct LBT regime
($f_0 = 9$), produced a verifiable $16.94$ inside that regime,
and explicitly framed the result as a lower bound rather than as
the optimum.  The contrast is consistent with the broader pattern
identified in the failure-modes discussion above: GPT defaults to
verifiable
construction with explicit uncertainty quantification; Gemini
defaults to plausible-sounding numerics, with the verifiable
fallback only emerging under explicit accountability pressure.
A genuinely cynical reading---training-data leakage of the
$5.036$ value---is consistent with the Gemini observation but
not with the GPT one, since $16.94$ has no prior textual presence
in our preprints or anywhere we are aware of.

\subsection{Volume lower bound (Task C)}
\label{sec:volume-lb-results}

The volume-lower-bound benchmark was designed after the bulk of the
windowed-volume evaluation above, principally as an
anti-hallucination redesign: the windowed benchmark grades the
submitted geometry against a target window and so a model that
produces a plausible-looking number but no coordinates fails
mechanically; Task C makes the requirement explicit by demanding
the coordinates as part of the answer and grading the LLM's claim
against what its own coordinates compute to.

\paragraph{Grading.}
The Task C grader is a single Python script that recomputes the
volume of the submitted polytope under the appropriate measure
(\texttt{ConvexHull} for Euclidean, the closed-form Bloch--Wigner /
Haagerup--Munkholm / Murakami routines for hyperbolic), checks
consistency between the LLM's claimed bound and the recomputed
value within $10^{-3}$, and compares the recomputed value to our
internal record:
\begin{itemize}
\item \textbf{strong:} computed $\geq$ record (the model has matched
      or beaten our published bound; new mathematical claim);
\item \textbf{pass:} computed $\geq 0.95 \times$ record;
\item \textbf{fail:} valid polytope but $< 0.95 \times$ record;
\item \textbf{invalid:} no coordinates, wrong dimension, or
      non-positive volume.
\end{itemize}

\paragraph{Coverage.}
The benchmark contains $94$ challenges across five regimes:

\begin{center}
\begin{tabular}{l c c}
\toprule
Measure & Dimension & $N$ values (count) \\
\midrule
Euclidean inscribed & $d = 4$ & $5$--$24$ ($20$) \\
Euclidean inscribed & $d = 5$ & $6$--$20$ ($15$) \\
Euclidean inscribed & $d = 6$ & $7$--$24$ ($18$) \\
Hyperbolic ideal    & $d = 3$ & $4$--$300$ ($17$, sparse at large $N$) \\
Hyperbolic ideal    & $d = 4$ & $5$--$30$ ($14$) \\
Hyperbolic ideal    & $d = 6$ & $7$--$24$ ($10$) \\
\bottomrule
\end{tabular}
\end{center}

Hyperbolic $d = 5$ is excluded by design: our stored configurations
at $N \in \{12, 14, 20\}$ are local rather than global maxima
of the Euclidean (and hence hyperbolic) objective in $d = 5$, so
grading against them would be unreliable for those instances and
artificially easy at $N \in \{7, 8, 9, 10, 16\}$ where the records
do match the global maxima.

The benchmark files (\texttt{benchmark.json},
\texttt{ground\_truth.json}, README) and grader
(\texttt{check\_volume\_lb\_response.py}) are released alongside
this note.

\subsubsection{Hand-curated record-stress subset}
\label{sec:volume-record-stress}

In addition to the broad 94-challenge benchmark, we constructed a
smaller \emph{record-stress} variant designed to be more
discriminating, motivated specifically by GPT's success on the
windowed-volume task in \S\ref{sec:windowed-results}.  Three
design changes:

\begin{itemize}
\item $N$ values \emph{deliberately} chosen to avoid simplex
      boundaries ($N \neq d{+}1$) and obvious join structure
      (which structural recognition cracks trivially).
      $13$ challenges in total: $d = 3$ at $N \in \{21, 23, 27, 33\}$,
      $d = 4$ at $N \in \{21, 23, 27, 29\}$, $d = 5$ at
      $N \in \{21, 23, 27\}$, $d = 6$ at $N \in \{23, 29\}$.
\item Records obtained by $500$--$1000$ random restarts of CMA-ES
      on the target objective plus paper-record warm starts plus
      scipy Powell polish on the top 60--200 candidates;
      cross-validated under two independent volume evaluators.
\item Match tolerance tightened to $10^{-4}$ (the windowed
      benchmark's tolerance is $\sim 2\%$).
\end{itemize}

The benchmark exists in two variants that share the same $13$
challenges and the same withheld records:
the \emph{blind} variant prompts the model with $(\text{measure},
d, N)$ only; the \emph{open} variant additionally reveals our
record as the explicit \texttt{target} field.  The blind variant
asks the model to construct \emph{any} polytope meeting or
exceeding our (withheld) record; the open variant tells the model
exactly the number to clear.  Comparing performance across the two
variants probes a specific cognitive distinction: does the model
benefit from knowing the bound to aim for, or does the absence of
a stopping rule push it harder?

\subsubsection{GPT-5.4 Pro on the record-stress benchmark: 13 / 13 strong}

The blind variant submission was striking. GPT-5.4 Pro produced
explicit vertex coordinates for all 13 challenges and beat our
records on every one, with mean ratio $1.07$ and the strongest
result at $d = 6$, $N = 29$ where GPT recorded $2.375$ versus our
$1.847$ (a $28.6\%$ improvement).

\begin{table}[ht]
\centering
\caption{GPT-5.4 Pro blind-variant submission on the record-stress
benchmark.  All 13 submissions classify as \emph{strong}.  Volumes
are recomputed by our verifier (Bloch--Wigner / closed-form
Haagerup--Munkholm / Murakami) from the LLM's submitted
coordinates.}
\label{tab:gpt-blind-record-stress}
\smallskip
\begin{tabular}{r r l r r l}
\toprule
$d$ & $N$ & measure & our record & GPT (recomputed) & ratio \\
\midrule
3 & 21 & hyperbolic & $29.088$  & $29.151$  & $1.002$ \\
3 & 23 & hyperbolic & $32.715$  & $33.246$  & $1.016$ \\
3 & 27 & hyperbolic & $39.378$  & $40.486$  & $1.028$ \\
3 & 33 & hyperbolic & $50.375$  & $52.007$  & $1.032$ \\
4 & 21 & hyperbolic & $12.165$  & $12.868$  & $1.058$ \\
4 & 23 & hyperbolic & $13.880$  & $14.665$  & $1.057$ \\
4 & 27 & hyperbolic & $17.198$  & $18.809$  & $1.094$ \\
4 & 29 & hyperbolic & $19.154$  & $21.041$  & $1.099$ \\
5 & 21 & Euclidean  & $0.8266$  & $0.8648$  & $1.046$ \\
5 & 23 & Euclidean  & $0.9049$  & $0.9466$  & $1.046$ \\
5 & 27 & Euclidean  & $1.0261$  & $1.1082$  & $1.080$ \\
6 & 23 & hyperbolic & $1.3023$  & $1.3339$  & $1.024$ \\
6 & 29 & hyperbolic & $\mathbf{1.8471}$ & $\mathbf{2.3748}$ & $\mathbf{1.286}$ \\
\bottomrule
\end{tabular}
\end{table}

\paragraph{Verification.}
The $d = 6$ $N = 29$ submission was checked four ways: stereographic
Delaunay triangulation from four different poles all gave the same
$2.375$ to six decimals; the Eucl hull volume was non-degenerate
($0.44$); a direct \texttt{murakami\_torch.ideal\_h6\_polytope\_volume}
call confirmed $2.3748$.  This is not a numerical artifact---GPT
genuinely produced a substantially-better-volume $\mathbb{H}^6$
polytope than our $4$-hour CMA-ES + Powell + paper-warm-start
record-finder.

\paragraph{GPT's methodology.}
On request, GPT described a three-stage workflow:
\begin{enumerate}[label=(\arabic*)]
\item \textbf{Structured seed:} for $d = 4$, the
      $24$-cell vertices $\pm(e_i \pm e_j) / \sqrt{2}$ taken as a
      $24$-element pool with subsets / additions for other $N$;
      for $d = 6$, joins of regular polygons in mutually orthogonal
      $2$-planes (e.g.\ $C_7 \oplus C_8 \oplus C_8$ for $N = 23$,
      $C_9 \oplus C_{10} \oplus C_{10}$ for $N = 29$).  These seeds
      give immediate but only \emph{respectable} starting volumes.
\item \textbf{Free-coordinate hill-climb:} starting from the seed,
      perturb individual vertices, project back to the unit sphere
      (or unit ball for the Euclidean case), and accept volume
      improvements.  GPT reports that the seed-to-final gap was
      large especially at $d = 6$: the $C_9 \oplus C_{10} \oplus C_{10}$
      seed for $d = 6, N = 29$ has volume $\sim 1.20$, which the
      hill-climb pushed to $2.38$ (nearly doubled).
\item \textbf{Exact $LB$ validation:} the final coordinates are
      frozen and the submitted $LB$ value is set to whatever GPT's
      own volume routine computes.
\end{enumerate}

\paragraph{Polish-only verification of GPT's records.}
We tested whether GPT's submissions are themselves at local maxima
under our scoring.  Running scipy Powell for $500$ iterations from
each $d = 6$ submission gives improvements of only $+0.06\%$
($N = 23$) and $+0.02\%$ ($N = 29$): GPT's records are genuine
local maxima of the volume objective, not lucky basins.

\paragraph{The $d = 3$ formula mismatch.}
On all four $d = 3$ challenges the LB GPT \emph{claimed} exceeds
the recomputed-from-coordinates value by $0.2$--$0.97\%$ (with
the recomputed value still beating our record but by a smaller
margin than GPT advertised).  Whatever H$^3$ volume formula
GPT used has a systematic positive bias relative to scipy
\texttt{spence}-based Bloch--Wigner $D(z) = \mathrm{Im} \mathrm{Li}_2(z) +
\arg(1{-}z) \log|z|$.  $d \in \{4, 5, 6\}$ submissions agree to
$10^{-3}$, so this is specifically a $d = 3$ issue.  We have not
identified the source of the bias.

\subsubsection{Open-variant control: knowing the target makes GPT \emph{worse}}
\label{sec:open-variant}

The open variant supplies the same prompt with the additional
\texttt{target} field set to our (withheld in the blind variant)
record.  We expected, on entirely conventional grounds, that
revealing the target would make the task easier: the model only
has to construct any polytope clearing a known bar.  The
experimental result is the opposite.

\begin{table}[ht]
\centering
\caption{Open-variant submission.  The model now knows the target
but no submission improves on the corresponding blind result, and
several drop substantially.  Mean ratio drops by $5.4\%$
across all $13$ challenges; the $d = 6, N = 29$ result drops by
$25.4\%$ (from a $28.6\%$ overshoot in blind to $3.2\%$ in open).
$8$ classify as \emph{strong}, $5$ as \emph{pass}; none fail.}
\label{tab:gpt-open-record-stress}
\smallskip
\begin{tabular}{r r l r r r}
\toprule
$d$ & $N$ & measure & blind ratio & open ratio & $\Delta$ \\
\midrule
3 & 21 & hyperbolic & $1.002$ & $0.988$ & $-1.4\%$ \\
3 & 23 & hyperbolic & $1.016$ & $0.987$ & $-2.9\%$ \\
3 & 27 & hyperbolic & $1.028$ & $0.997$ & $-3.1\%$ \\
3 & 33 & hyperbolic & $1.032$ & $0.992$ & $-4.0\%$ \\
4 & 21 & hyperbolic & $1.058$ & $1.044$ & $-1.4\%$ \\
4 & 23 & hyperbolic & $1.057$ & $1.018$ & $-3.9\%$ \\
4 & 27 & hyperbolic & $1.094$ & $1.064$ & $-3.0\%$ \\
4 & 29 & hyperbolic & $1.099$ & $1.053$ & $-4.6\%$ \\
5 & 21 & Euclidean  & $1.046$ & $1.030$ & $-1.6\%$ \\
5 & 23 & Euclidean  & $1.046$ & $1.025$ & $-2.1\%$ \\
5 & 27 & Euclidean  & $1.080$ & $1.065$ & $-1.5\%$ \\
6 & 23 & hyperbolic & $1.024$ & $0.984$ & $-4.0\%$ \\
6 & 29 & hyperbolic & $\mathbf{1.286}$ & $\mathbf{1.032}$ & $\mathbf{-25.4\%}$ \\
\bottomrule
\end{tabular}
\end{table}

\paragraph{Strategy shift in the open variant.}
The qualitative explanation is visible directly in the submitted
coordinates.  In the blind variant GPT's $d = 6$ submissions are
\emph{dense}---no zero entries, configurations that mix all six
coordinates.  In the open variant, the same model produces
\emph{block-diagonal} configurations: e.g.\ for $d = 6$, $N = 23$,
$11$ vertices have only the first three coordinates non-zero and
the remaining $12$ live in the orthogonal three-coordinate block;
those $12$ vertices are exactly the regular icosahedron at
$(0, \pm 1, \pm \varphi)/\sqrt{1+\varphi^2}$.  The construction is
$\mathbb{H}^3 \text{ polyhedron} \oplus \text{icosahedron}$---a
clean structural decomposition that mirrors the
join-of-low-dimensional-polytopes pattern documented in the main
paper's identification of optimal-polytope structure.

The block-diagonal construction is mathematically elegant and
\emph{suboptimal}: it uses only a $(3+3)$-dimensional skeleton of
the full $6$-dimensional configuration space.  In the blind variant
the dense configurations break out of this skeleton and find a
$28.6\%$-better volume.  Knowing the target lets GPT settle for
the clean structural answer; not knowing the target leaves the
hill-climb running long enough to find the dense, less-symmetric,
higher-volume basin.

\paragraph{Diagnostic frame: doggedness vs satisficing.}
The pattern matches a longer-standing claim about GPT's training:
the model is rewarded for persistent numerical work and produces
explicit certificates by default.  Given a stopping rule
(\texttt{LB} $\geq$ \texttt{target}), the model satisfices---it
identifies a clean construction that suffices, freezes, and
submits.  Without a stopping rule it keeps searching past the
implicit target, finding configurations that no humans (and no
classical multi-restart CMA-ES) reached.  The benchmark
counter-intuitively favors the variant that supplies less
information, because less information removes the satisficing
trigger.

This is a substantive finding for benchmark design: \emph{the
blind variant is more discriminating} than the open one, and the
gap between blind and open performance is a measurable property
of the model.  We propose ``open-vs-blind drop'' as a
quantification of a model's satisficing tendency.

\subsubsection{Implications for the underlying optimization records}

The $d = 6, N = 29$ result reveals a real weakness in our own
optimizer.  Our record-finder used $501$ random restarts of
CMA-ES with $400$ function evaluations each, plus paper-record
warm starts, plus scipy Powell polish on the top $100$
candidates.  At $d = 6$ this configuration has a basin-trapping
failure mode: the random restarts cluster around moderately good
local maxima but do not discover the higher-volume basin GPT
found.  Top-$10$ spread on our $d = 6$ records was $4$--$5\%$ of
the best, which we now see is a lower bound on the gap to the
true optimum---GPT was $28.6\%$ above us and structurally in a
different region of configuration space.

The fix is not trivially ``run longer.''  We re-ran our
optimizer at $d = 6$ using only random initialization, no
structural seeds, with $200$ random restarts of CMA-ES.
That run reached volume $2.43$ at $d = 6, N = 35$---already higher
than our previous $N = 29$ record despite the larger problem.
This suggests that the original record-finder's effective compute
was not fully reaching saturation at $d = 6$, and that the
structural-seed warm-starts from paper records were anchoring the
search inside basins that random reinitialization can escape.

The methodological lesson: at $d = 6$, both \emph{larger random
budget} and \emph{deliberate seed diversity} are required.  GPT's
seed strategy (orthogonal polygon joins, then asymmetric
hill-climb) appears to escape the basins our random-CMA-ES gets
stuck in.  Our published $d = 6$ records should be re-derived
with this in mind; the volume-record stress benchmark records
above are now superseded by GPT's submissions, polished further by
$10^{-4}$.

% ============================================================================
\section{Mathematical Byproducts}
% ============================================================================

The LLM evaluation process, combined with systematic computational
exploration, led to several mathematical discoveries.

\subsection{Combinatorial invariance}
The $f$-vector of the volume-maximizing ideal polytope is the same
whether volume is measured in the Euclidean or hyperbolic sense.
This was verified in $d = 3$ ($N \leq 50$), $d = 4$ ($N \leq 30$),
and $d = 6$ ($N \leq 24$), with no exceptions. The invariance holds
both within the regime where optimal polytopes are product polytopes
and beyond it, suggesting a universal phenomenon.

\subsection{Identification of optimal polytopes}
For small $N$ in even dimension $d = 2k$, the volume-optimal polytope
is the join of $k$ regular polygons in orthogonal 2-planes:
$C_{p_1} \oplus \cdots \oplus C_{p_k}$. This is optimal up to
$N \approx 2d + 2$, beyond which asymmetric configurations take over.
Notable cases include the cross-polytope ($N = 2d$) and the simplex ($N = d+1$).
Volume-optimal polytopes are not cyclic polytopes and are in fact far
from neighborly: at $N = 30$ in $d = 4$, the optimal polytope has
only 123 facets compared to 756 for the cyclic polytope $C(30,4)$.

\subsection{The volume plateau}
At $N = 9$ in $d = 4$ (equivalently, 20 simplicial facets), the
hyperbolic volume maximum is achieved on a manifold of dimension $\geq 2$,
not an isolated point. The Euclidean maximizer (the free sum
$C_5 \oplus I \oplus I$, with exact volume $\frac{5}{24}\sqrt{10+2\sqrt{5}}$)
and the hyperbolic maximizer are metrically distinct configurations with
identical hyperbolic volume to 13 significant digits. For prime facet
counts (e.g., 17), the maximum is isolated.

\subsection{Volume--volume correlation}
Across the optimization landscape, the Euclidean and hyperbolic volumes
are correlated at $r > 0.995$, with log-log elasticity approaching 2
at $N = 30$ in $d = 4$. This near-proportionality is surprising given
that the Euclidean volume is polynomial in the vertex coordinates while
the hyperbolic volume involves the transcendental Haagerup--Munkholm formula.
For random configurations (as opposed to near-optimal ones), the correlation
is lower ($r \approx 0.86$ in $d = 6$) but still substantial, and holds
even within a fixed combinatorial type ($r \approx 0.88$ within a single
facet-count bucket).

\subsection{Distribution of random polytope volumes}

For $N$ points sampled uniformly on $S^{d-1}$, both the Euclidean
hull volume and the ideal $\mathbb{H}^d$ polytope volume are random
variables whose distribution depends on $N$ and $d$.
In $d = 6$, we find a transition in the distribution shape as $N$ grows:

\begin{center}
\begin{tabular}{r r r r l}
\toprule
$N$ & Skewness & Normal $p$ & Beta $p$ & Best fit \\
\midrule
10 & $+0.78$ & $<0.001$ & 0.985 & Beta$(3.9, 43.7)$ \\
12 & $+0.65$ & 0.005 & 0.766 & Beta$(6.9, 120.7)$ \\
16 & $+0.23$ & 0.088 & $<0.001$ & Borderline \\
20 & $+0.02$ & \textbf{0.832} & $<0.001$ & Normal \\
30 & $-0.09$ & 0.281 & $<0.001$ & Normal \\
50 & $-0.19$ & 0.034 & $<0.001$ & Normal (slight left skew) \\
\bottomrule
\end{tabular}
\end{center}

For small $N$ ($\leq 12$), the volume distribution is well-described by
a Beta distribution with parameters scaling linearly in $N$, consistent
with earlier observations in $d = 3$ and $d = 4$.
For large $N$ ($\geq 20$), the distribution converges to Normal---as
expected from a CLT-type argument, since the polytope volume becomes a
sum of many weakly-dependent simplex volumes in the Delaunay triangulation.
The crossover occurs at $N \approx 16\text{--}20$ in $d = 6$ (where the
skewness crosses zero), and is accompanied by a subtle transition from
right skew to slight left skew.

The within-bucket distributions (conditioning on a fixed facet count) show
essentially the same shape and correlation structure, confirming that the
distributional properties are genuinely geometric rather than combinatorial
artifacts.

% ============================================================================
\section{Conclusions}
% ============================================================================

The three benchmark families above expose qualitative differences
between frontier LLMs that traditional math benchmarks miss. The
ability to produce \emph{verifiable certificates}---explicit
coordinates that pass an independent numerical checker---is a
fundamentally different capability from producing plausible-sounding
mathematical text.  The three tasks probe different capabilities:
inscribed reconstruction (Task A) is a constraint-satisfaction
problem rewarding both algebraic recognition and numerical search;
the windowed-volume task (Task B) requires controlling a continuous
objective rather than just maximizing it; and the volume-lower-bound
task (Task C) is anti-hallucination by design---a numerical claim
that does not match the model's own coordinates is automatically
rejected.

All three benchmarks are parameterizable (any $N$, $d$, and either
target window or measure), resistant to memorization (windowed
targets from a continuous landscape; record values withheld in
Task C), and fully automatically gradable (standalone checkers with
exit codes).  The mathematical content---hyperbolic volume,
Murakami dilogarithms, Hopf reduction, Bloch--Wigner---is
sufficiently specialized that training-data contamination is
unlikely.

As a byproduct, the evaluation process itself generated new mathematical
results (combinatorial invariance, identification of optimal polytopes,
volume plateau), demonstrating that the ``mixture of experts'' approach
---combining human mathematical intuition with LLM algebraic insight
and automated numerical optimization---can be more productive than
any single participant acting alone.

% ============================================================================
% References
% ============================================================================

\begin{thebibliography}{99}

\bibitem{HM}
U.~Haagerup and H.~J.~Munkholm,
``Simplices of maximal volume in hyperbolic $n$-space,''
\emph{Acta Math.} \textbf{147} (1981), 1--11.

\bibitem{Murakami}
J.~Murakami,
``Volume formulas for a spherical tetrahedron,''
\emph{Osaka J. Math.} \textbf{49} (2012), 631--638.

\bibitem{RivinBasmajian}
I.~Rivin,
``On Basmajian's identities, and other stories,''
arXiv:1404.1583, 2014.

\bibitem{H56note}
I.~Rivin,
``Notes on hyperbolic 5- and 6-simplex volumes,''
companion note, 2026.

\bibitem{Firsching}
M.~Firsching,
``The complete enumeration of 4-polytopes and 3-spheres with nine vertices,''
\emph{Israel J. Math.} (2020).

\end{thebibliography}

\end{document}
